\usepackage{verbatim, multicol, tabularx,hyperref, tikz, enumitem}
\usepackage{amsmath,amsthm, amssymb, stmaryrd, latexsym, bm, listings, qtree}
\usepackage[margin=1in]{geometry}

\lstset{
  extendedchars=\true,
  inputencoding=utf8,
  literate=
  {é}{{\'{e}}}1
  {è}{{\`{e}}}1
  {ê}{{\^{e}}}1
  {ë}{{\¨{e}}}1
  {û}{{\^{u}}}1
  {ù}{{\`{u}}}1
  {â}{{\^{a}}}1
  {à}{{\`{a}}}1
  {î}{{\^{i}}}1
  {ô}{{\^{o}}}1
  {ç}{{\c{c}}}1
  {Ç}{{\c{C}}}1
  {É}{{\'{E}}}1
  {Ê}{{\^{E}}}1
  {À}{{\`{A}}}1
  {Â}{{\^{A}}}1
  {Î}{{\^{I}}}1
  {Ö}{{\"O}}1
  {Ä}{{\"A}}1
  {Ü}{{\"U}}1
  {ö}{{\"o}}1
  {ä}{{\"a}}1
  {ü}{{\"u}}1
  {ß}{{\ss}}1
  ,
  aboveskip=1mm,
  belowskip=1mm,
  showstringspaces=false,
  columns=flexible,
  basicstyle={\scriptsize\ttfamily},
  numbers=left,
  frame=single,
  framextopmargin=0pt,
  framexbottommargin=0pt,
  breaklines=true,
  breakatwhitespace=true,
  keywordstyle=\color{blue},
  identifierstyle=\color{violet},
  stringstyle=\color{teal},
  commentstyle=\color{darkgray}
}

\hypersetup{colorlinks=true,urlcolor=blue}

\headheight = 0.05 in
\headsep = 0.05 in
\parskip = 0.05in
\parindent = 0.0in
\floatsep = 0.05in

\DeclareMathOperator*{\argmin}{arg\!min}
\DeclareMathOperator*{\argmax}{arg\!max}

\title{Introduction to AI Review}
\author{Artificial Intelligence}
\date{}

\begin{document}
\maketitle

\begin{questions}


\setcounter{section}{0} % So that next \section is 1


\question What is intelligence?

\begin{solution}[1.5in]
Intelligence is the ability to do:

\begin{itemize}
\item {\bf Problem solving}: transforming a given state into a goal state.
\item {\bf Inference}: deriving new knowledge from existing knowledge.
\item {\bf Reasoning}: answering questions or justifying answers.
\item {\bf Decision making}: choosing actions that maximize expected utility.
\item {\bf Learning}: using experience (e.g., data) to get better at a task, such as those listed above.
\end{itemize}

\end{solution}

\question What is artificial intelligence?

\begin{solution}[1.5in]
A human-made (artificial) system that does any of the above is an AI system.
\end{solution}

\question What are the four approaches to AI?

\begin{solution}[2in]
\end{solution}

\newpage

\question Explain ``thinking humanly.''

\begin{solution}[1.5in]
\end{solution}

\question Explain ``acting humanly.''

\begin{solution}[1.5in]
\end{solution}

\question What is the Turing Test?

\begin{solution}[1.5in]
A computer is intelligent if a human interrogator, after posing written questions, can’t tell whether written responses come from a person or from the computer.
\end{solution}

\question Explain ``thinking rationally.''

\begin{solution}[1.5in]
\end{solution}

\question Explain ``acting rationally.''

\begin{solution}[1.5in]
\end{solution}

\newpage

\question When was the term ``Artificial Intelligence'' coined, and by whom?

\begin{solution}[1in]
John McCarthy coined the term artificial intelligence in 1955 in his proposal for the first workshop on AI to take place at Dartmouth College in the summer of 1956.
\end{solution}

\question Geoff Hinton has been called the ``godfather of AI.''  How old was Geoff Hinton when the field of AI was launched?

\begin{solution}[1in]
8
\end{solution}

\question What is the {\bf strong AI} hypothesis?

\begin{solution}[1in]
The assertion that machines that execute intelligent behavior are actually consciously thinking.
\end{solution}

\question What is the {\bf weak AI} hypothesis?

\begin{solution}[1in]
The idea that machines can only {\it act} as if they are intelligent, that is, they can {\it simulate} intelligent behavior.
\end{solution}

\question What is {\bf narrow AI}?

\begin{solution}[1in]
Algorithms and models for accomplishing specific tasks, such as playing chess or recognizing objects in an image.
\end{solution}

\question What is {\bf general AI}?

\begin{solution}[1in]
Algorithms and models for accomplishing any intelligent task.  Many consider this to be equivalent to ``human-leval AI,'' but in general AGI (artificial general intelligence) is ill-defined and mostly a marketing buzzword.
\end{solution}

\newpage

\question Is artificial intelligence possible?

\begin{solution}[2in]
The answer depends what kind of AI you mean.  Narrow AI is of course possible.  General AI, assuming it were rigorously defined at all, may or may not be possible.  Weak AI is, of course, possible by construction -- any narrow AI system is an example of weak AI.  Strong AI is mostly a philosophical question.
\end{solution}


\end{questions}

\end{document}
